\documentclass[a4paper,12pt]{article}

\usepackage[utf8]{inputenc}
\usepackage[T1]{fontenc}
\usepackage[italian]{babel}
\usepackage{amsmath, amssymb}
\usepackage{graphicx}
\usepackage{hyperref}
\usepackage{geometry}
\usepackage{booktabs}
\usepackage{float}

\geometry{margin=2.5cm}

\title{Report 2}

\begin{document}

\maketitle

\newpage

\section{Disclaimer}
Il seguente report è da considerarsi un' estensione del file \emph{Sviluppo\_Classical\_ksp.pdf}. 
Come test sono stati usati ancora \texttt{ksp\_1.txt} e \texttt{ksp\_2.txt} . 

\section{Introduzione}\label{sec:Introduzione}
Nel seguente documento sono esposti i risultati ottenuti con i seguenti cambiamenti:

\begin{itemize}
    \item Aggiunto un nuovo metodo per il calcolo di
    \[
    \lambda = \frac{\sum_{i=1}^{n} p_i}{650 \cdot C}
    \]

    \item Modificato il parametro \texttt{k} nelle funzioni \texttt{script.annealer()} e \texttt{solver.solve()}.
    
    \item Aggiunto test con un altro tipo di ottimizzatore: \texttt{Gurobi}.
    
    \item Modifica come e quali dati stampare in output.
\end{itemize}

\subsection{Soluzioni trovate con Gurobi}
Su test case \texttt{ksp\_1.txt}, considerando i due valori di
\[
\lambda = \frac{\sum_{i=1}^{n} p_i}{650 \cdot C}
\qquad \text{e} \qquad
\lambda = \frac{\sum_{i=1}^{n} p_i}{3},
\]
si ottengono i seguenti risultati:

\begin{table}[H]
\centering
\begin{tabular}{lc}
\toprule
Numero totale di run & 1000 \\
Best Weight & 101 \\
Best Profit & 198 \\
Gurobi Weight & 101 \\
Gurobi Profit & 198 \\
fQ & -1455540.67 \\
Items & \texttt{[1, 1, 0, 0, 0, 0, 1, 0, 0, 0]} \\
\bottomrule
\end{tabular}
\caption{Risultati con $\lambda = \displaystyle \frac{\sum_{i=1}^{n} p_i}{3}$}
\label{tab:tab1}
\end{table}

\begin{table}[H]
\centering
\begin{tabular}{lc}
\toprule
Numero totale di run & 1000 \\
Best Weight & 101 \\
Best Profit & 198 \\
Gurobi Weight & 177 \\
Gurobi Profit & 302 \\
fQ & -330.85 \\
Items & \texttt{[1, 1, 0, 0, 0, 1, 1, 1, 0, 0]} \\
\bottomrule
\end{tabular}
\caption{Risultati con $\lambda = \displaystyle \frac{\sum_{i=1}^{n} p_i}{650 \cdot C}$}
\label{tab:tab2}
\end{table}

In questo caso Gurobi sembra trovare la soluzione migliore solo nel caso si utilizzi $\lambda = \displaystyle \frac{\sum_{i=1}^{n} p_i}{3}$.

Su test case \texttt{ksp\_2.txt} invece:

\begin{table}[H]
\centering
\begin{tabular}{lc}
\toprule
Numero totale di run & 1000 \\
Best Weight & 978 \\
Best Profit & 7740 \\
Gurobi Weight & 1001 \\
Gurobi Profit & 6288 \\
fQ & -10200710468.33 \\
\bottomrule
\end{tabular}
\caption{Risultati con $\lambda = \displaystyle \frac{\sum_{i=1}^{n} p_i}{3}$}
\label{tab:tab3}
\end{table}

\begin{table}[H]
\centering
\begin{tabular}{lc}
\toprule
Numero totale di run & 1000 \\
Best Weight & 978 \\
Best Profit & 7740 \\
Gurobi Weight & 978 \\
Gurobi Profit & 7740 \\
fQ & -54748.31 \\
\bottomrule
\end{tabular}
\caption{Risultati con $\lambda = \displaystyle \frac{\sum_{i=1}^{n} p_i}{650 \cdot C}$}
\label{tab:tab4}
\end{table}

Nel secondo test case invece è il contrario: trova la soluzione con $\lambda = \displaystyle \frac{\sum_{i=1}^{n} p_i}{650 \cdot C}$.

\section{Soluzioni trovate con l'annealer}
Negli esempi, come anticipato in Sezione~\ref{sec:Introduzione}, ho cambiato il parametro \texttt{k = 10}.

\subsection{NON utilizzando QALS}

Per \texttt{ksp\_1.txt} si ha:

\begin{table}[H]
\centering
\begin{tabular}{lc}
\toprule
Numero totale di run & 1000 \\
Best Weight & 101 \\
Best Profit & 198 \\
Avg Weight GAP & -0.5 \\
Avg Profit GAP & 67.2 \\
Avg fQ & -1455294.95 \\
fQ Min Fuound & -1455540.67 \\
n of fQ Min Found & 424 \\
Items & \texttt{[1, 1, 0, 0, 0, 0, 1, 0, 0, 0]} \\
\bottomrule
\end{tabular}
\caption{Risultati con $\lambda = \displaystyle \frac{\sum_{i=1}^{n} p_i}{3}$ \ref{tab:tab1}}
\end{table}

Nel primo test pare l'annealer riesca a trovare la soluzione corretta con una percentuale del 42\% (il 
confronto è fatto con Gurobi/soluzione ottima).

\begin{table}[H]
\centering
\begin{tabular}{lc}
\toprule
Numero totale di run & 1000 \\
Best Weight & 101 \\
Best Profit & 198 \\
Avg Weight GAP & -76.0 \\
Avg Profit GAP & -104.0 \\
Avg fQ & -330.85 \\
fQ Min Fuound & -330.85 \\
n of fQ Min Found & 1000 \\
Items & \texttt{[1, 1, 0, 0, 0, 1, 1, 1, 0, 0]} \\
\bottomrule
\end{tabular}
\caption{Risultati con $\lambda = \displaystyle \frac{\sum_{i=1}^{n} p_i}{650 \cdot C}$ \ref{tab:tab2}}
\end{table}

Nel secondo invece le solzioni trovate dall'ennealer sono tutte uguali tra loro e a quelle trovate 
da Gurobi.

Per quanto riguarda \texttt{ksp\_2.txt}:

\begin{table}[H]
\centering
\begin{tabular}{lc}
\toprule
Numero totale di run & 1000 \\
Best Weight & 978 \\
Best Profit & 7740 \\
Avg Weight GAP & -23.0 \\
Avg Profit GAP & 3885.4 \\
Avg fQ & -10200678481.45 \\
fQ Min Fuound & -10200710468.33 \\
n of fQ Min Found & 1 \\
\bottomrule
\end{tabular}
\caption{Risultati con $\lambda = \displaystyle \frac{\sum_{i=1}^{n} p_i}{3}$ \ref{tab:tab3}}
\end{table}

L'unica solzione trovata (non ottima) anche in questo caso è la stessa di quella trovata da
Gurobi.

\begin{table}[H]
\centering
\begin{tabular}{lc}
\toprule
Numero totale di run & 1000 \\
Best Weight & 978 \\
Best Profit & 7740 \\
Avg Weight GAP & -32.1 \\
Avg Profit GAP & 72.6 \\
Avg fQ & -54603.47 \\
fQ Min Fuound & -54748.31 \\
n of fQ Min Found & 617 \\
\bottomrule
\end{tabular}
\caption{Risultati con $\lambda = \displaystyle \frac{\sum_{i=1}^{n} p_i}{650 \cdot C}$ \ref{tab:tab4}}
\end{table}

In questo le soluzioni ottime trovate sono il 61\% e corrispondono a quelle trovate anche da Gurobi.

\subsection{Utilizzando QALS}

\texttt{ksp\_1.txt}:

\begin{table}[H]
\centering
\begin{tabular}{lc}
\toprule
Numero totale di run & 1000 \\
Best Weight & 101 \\
Best Profit & 198 \\
Avg Weight GAP & -75.1 \\
Avg Profit GAP & 22.7 \\
Avg fQ & -332380.35 \\
fQ Min Fuound & -1455540.67 \\
n of fQ Min Found & 2 \\
Items & \texttt{[1, 1, 0, 0, 0, 0, 1, 0, 0, 0]} \\
\bottomrule
\end{tabular}
\caption{Risultati con $\lambda = \displaystyle \frac{\sum_{i=1}^{n} p_i}{3}$ \ref{tab:tab1}}
\end{table}

Per questo caso di $\lambda$\ QALS riesce a trovare la stessa soluzione di Gurobi e quella ottima.

\begin{table}[H]
\centering
\begin{tabular}{lc}
\toprule
Numero totale di run & 1000 \\
Best Weight & 101 \\
Best Profit & 198 \\
Avg Weight GAP & -41 \\
Avg Profit GAP & 60.8 \\
Avg fQ & -126.04 \\
fQ Min Fuound & -284.04 \\
n of fQ Min Found & 3 \\
Items & \texttt{[1, 0, 0, 0, 0, 1, 1, 0, 0, 0]} \\
\bottomrule
\end{tabular}
\caption{Risultati con $\lambda = \displaystyle \frac{\sum_{i=1}^{n} p_i}{650 \cdot C}$}
\end{table}

Per il secondo valore di $\lambda$\ invece sembra ancora non riuscire a trovare la soluzione migliore
ma alcune con dei GAP non trascurabili.


\texttt{ksp\_2.txt}:

\begin{table}[H]
\centering
\begin{tabular}{lc}
\toprule
Numero totale di run & 1000 \\
Best Weight & 978 \\
Best Profit & 7740 \\
Avg Weight GAP & -8703.8 \\
Avg Profit GAP & -5790.5 \\
Avg fQ & 770684171363.27 \\
fQ Min Fuound & 170922927950.67 \\
n of fQ Min Found & 1 \\
\bottomrule
\end{tabular}
\caption{Risultati con $\lambda = \displaystyle \frac{\sum_{i=1}^{n} p_i}{3}$ \ref{tab:tab3}}
\end{table}

\begin{table}[H]
\centering
\begin{tabular}{lc}
\toprule
Numero totale di run & 1000 \\
Best Weight & 978 \\
Best Profit & 7740 \\
Avg Weight GAP & -8433.1 \\
Avg Profit GAP & -5337.1 \\
Avg fQ & 3340506.07 \\
fQ Min Fuound & 961014.64 \\
n of fQ Min Found & 1 \\
\bottomrule
\end{tabular}
\caption{Risultati con $\lambda = \displaystyle \frac{\sum_{i=1}^{n} p_i}{650 \cdot C}$ \ref{tab:tab4}}
\end{table}

In entrambi i casi le soluzioni trovate sono ben distanti sia da quella ottima che da quella trovata da Gurobi.

\end{document}